\documentclass[11pt, a4paper]{article}

% --- Packages ---
\usepackage[a4paper, top=2.5cm, bottom=2.5cm, left=2.5cm, right=2.5cm]{geometry}
\usepackage{amsmath, amssymb}
\usepackage{graphicx}
\usepackage{booktabs}
\usepackage{parskip}
\usepackage{microtype}
\usepackage{fancyhdr}
\usepackage{titlesec}
\usepackage{xcolor}
\usepackage{mdframed}
\usepackage{enumitem}
\usepackage{hyperref}
\usepackage{caption}
\usepackage{float}
\usepackage{array}

\hypersetup{colorlinks=true, linkcolor=black, urlcolor=blue, citecolor=black}

% --- Page style ---
\pagestyle{fancy}
\fancyhf{}
\renewcommand{\headrulewidth}{0.4pt}
\renewcommand{\footrulewidth}{0.4pt}
\fancyhead[L]{\small\textit{Study Notes: C-MAPSS Dataset}}
\fancyhead[R]{\small\thepage}
\fancyfoot[C]{\small\textit{``I'm not stupid, I'm just slow''}}

% --- Section style ---
\titleformat{\section}{\large\bfseries}{}{0em}{}[\titlerule]
\titleformat{\subsection}{\normalsize\bfseries}{}{0em}{}
\titleformat{\subsubsection}{\normalsize\itshape}{}{0em}{}

% --- Highlight boxes ---
\newmdenv[
  backgroundcolor=gray!8,
  linecolor=gray!40,
  linewidth=0.8pt,
  innerleftmargin=10pt,
  innerrightmargin=10pt,
  innertopmargin=8pt,
  innerbottommargin=8pt,
  skipabove=8pt,
  skipbelow=8pt
]{keybox}

\newmdenv[
  backgroundcolor=orange!6,
  linecolor=orange!40,
  linewidth=0.8pt,
  innerleftmargin=10pt,
  innerrightmargin=10pt,
  innertopmargin=8pt,
  innerbottommargin=8pt,
  skipabove=8pt,
  skipbelow=8pt
]{insightbox}

% --- Caption style ---
\captionsetup{font=small, labelfont=bf, skip=4pt}

% -------------------------------------------------------------------
\begin{document}
% -------------------------------------------------------------------

{\Large\bfseries Dataset Notes: C-MAPSS\\
Turbofan Engine Degradation Simulation Data}

\smallskip
{\small Saxena, Goebel, Simon, Eklund \quad --- \quad PHM'08 dataset \quad --- \quad
\textit{Working reference edition}}

\medskip\hrule

\tableofcontents

\bigskip\hrule\bigskip

% -------------------------------------------------------------------
\section{Overview}
% -------------------------------------------------------------------

The C-MAPSS dataset is the data product of the paper
\textit{Damage Propagation Modeling for Aircraft Engine
Run-to-Failure Simulation} (Saxena et al., PHM'08).
It is one of the most widely used benchmarks in predictive
maintenance and Remaining Useful Life (RUL) research.

This document is a working reference for the dataset files
themselves --- their structure, content, and how to use them.
It is intentionally separate from the paper notes. Understanding
the simulation methodology is one thing; working with the
actual files is another.

\begin{keybox}
\textbf{What the dataset is:}
A collection of multivariate time series, each representing
one simulated aircraft engine's sensor readings from initial
operation through failure. The task is to predict how many
cycles remain before failure, given only the sensor readings
up to the current cycle.
\end{keybox}

% -------------------------------------------------------------------
\section{The Four Sub-Datasets}
% -------------------------------------------------------------------

The C-MAPSS release is not one dataset but four, named
FD001 through FD004. They share the same 26-column file
format and the same engine model, but differ along two
independent axes of difficulty.

Figure~\ref{fig:matrix} shows the complete overview.

\begin{figure}[H]
  \centering
  \includegraphics[width=\textwidth]{09_dataset_matrix.png}
  \caption{The four C-MAPSS sub-datasets. Color encodes
  difficulty: teal = easiest, amber = medium, coral = hardest.
  The difficulty arrow on the right shows the overall ordering.}
  \label{fig:matrix}
\end{figure}

\subsection{Two independent axes of difficulty}

\textbf{Conditions axis} separates FD001/FD003 (1 operating
condition: sea level) from FD002/FD004 (6 operating conditions:
combinations of altitude, Mach number, and TRA). With multiple
conditions, raw sensor values are not comparable across cycles
because the same sensor reads differently at different flight
points. This requires normalization before any analysis.

\textbf{Fault mode axis} separates FD001/FD002 (HPC
degradation only) from FD003/FD004 (HPC degradation plus Fan
degradation). With two fault modes, the algorithm does not
know which component is degrading. No label is provided.

\begin{center}
\begin{tabular}{@{}lcccc@{}}
  \toprule
  & FD001 & FD002 & FD003 & FD004 \\
  \midrule
  Conditions  & 1 & 6 & 1 & 6 \\
  Fault modes & 1 & 1 & 2 & 2 \\
  Train       & 100 & 260 & 100 & 248 \\
  Test        & 100 & 259 & 100 & 249 \\
  Difficulty  & lowest & medium & medium & highest \\
  \bottomrule
\end{tabular}
\end{center}

\begin{insightbox}
\textbf{FD002 vs FD003 are incomparable.} Both are ``medium''
difficulty but hard in different ways. FD002 is hard because
of operating condition variability. FD003 is hard because of
fault mode ambiguity. An algorithm that handles FD002 well
may still fail on FD003, and vice versa.
\end{insightbox}

\begin{keybox}
\textbf{Practical warning:} Most published results report
on FD001 only --- it is the easiest, the most reproducible,
and the most cited. A result on FD001 alone should be
interpreted cautiously. It says nothing about robustness to
condition variability or fault mode ambiguity.
\end{keybox}

% -------------------------------------------------------------------
\section{File Structure: 12 Files}
% -------------------------------------------------------------------

Each sub-dataset produces exactly three files, giving
$4 \times 3 = \mathbf{12}$ files in total:

\begin{center}
\begin{tabular}{@{}lp{10cm}@{}}
  \toprule
  File & Contents \\
  \midrule
  \texttt{train\_FDxxx.txt} &
    Full run-to-failure trajectories. Each engine runs from
    its initial condition to failure. The last row of each
    engine is the failure cycle.
    \textbf{RUL is not provided} --- it must be computed
    from the data: $\text{RUL}(t) = t_{\max} - t$. \\[6pt]
  \texttt{test\_FDxxx.txt}  &
    Truncated trajectories. Each engine's time series stops
    at some point before failure. The last row is
    \emph{not} the failure cycle.
    \textbf{RUL is unknown} and must be predicted. \\[6pt]
  \texttt{RUL\_FDxxx.txt}   &
    Ground truth for the test set. One integer per engine:
    the true RUL at the last observed cycle of that engine.
    Used only for evaluation. \\
  \bottomrule
\end{tabular}
\end{center}

Figure~\ref{fig:filerel} shows how the three files relate to
a single conceptual engine life.

\begin{figure}[H]
  \centering
  \includegraphics[width=0.92\textwidth]{10_train_test_rul_relationship.png}
  \caption{The relationship between the three file types for
  one engine. The train file covers the full life to failure.
  The test file stops at $t^*$. The RUL file provides one
  number --- the true cycles remaining at $t^*$.}
  \label{fig:filerel}
\end{figure}

All \texttt{train} and \texttt{test} files are plain text,
space-separated, with no header row. The \texttt{RUL} files
are a single column of integers, one per engine, also with
no header.

% -------------------------------------------------------------------
\section{The 26-Column Format}
% -------------------------------------------------------------------

Every row in a \texttt{train} or \texttt{test} file is one
cycle of one engine. The 26 columns are fixed across all
four sub-datasets.

\begin{table}[H]
  \centering
  \caption{Complete column reference for C-MAPSS train and
  test files. This table is your Excel header row.}
  \label{tab:columns}
  \small
  \begin{tabular}{@{}rllp{6cm}@{}}
    \toprule
    Col & Name & Units & Description \\
    \midrule
    1  & \texttt{unit\_number}   & ---     & Engine ID. Integer, resets per file. \\
    2  & \texttt{time\_cycles}   & cycles  & Cycle counter per engine. Starts at 1. \\
    \midrule
    3  & \texttt{op\_setting\_1} & ft      & Altitude. \\
    4  & \texttt{op\_setting\_2} & ---     & Mach number. \\
    5  & \texttt{op\_setting\_3} & ---     & Throttle resolver angle (TRA). \\
    \midrule
    6  & \texttt{T2}             & \textdegree R & Total temperature, fan inlet. \\
    7  & \texttt{T24}            & \textdegree R & Total temperature, LPC outlet. \\
    8  & \texttt{T30}            & \textdegree R & Total temperature, HPC outlet. \\
    9  & \texttt{T50}            & \textdegree R & Total temperature, LPT outlet. \\
    10 & \texttt{P2}             & psia    & Pressure, fan inlet. \\
    11 & \texttt{P15}            & psia    & Total pressure, bypass duct. \\
    12 & \texttt{P30}            & psia    & Total pressure, HPC outlet. \\
    13 & \texttt{Nf}             & rpm     & Physical fan speed. \\
    14 & \texttt{Nc}             & rpm     & Physical core speed. \\
    15 & \texttt{epr}            & ---     & Engine pressure ratio ($P_{50}/P_2$). \\
    16 & \texttt{Ps30}           & psia    & Static pressure, HPC outlet. \\
    17 & \texttt{phi}            & pps/psi & Fuel flow / $P_{s30}$ ratio. \\
    18 & \texttt{NRf}            & rpm     & Corrected fan speed. \\
    19 & \texttt{NRc}            & rpm     & Corrected core speed. \\
    20 & \texttt{BPR}            & ---     & Bypass ratio. \\
    21 & \texttt{farB}           & ---     & Burner fuel-air ratio. \\
    22 & \texttt{htBleed}        & ---     & Bleed enthalpy. \\
    23 & \texttt{Nf\_dmd}        & rpm     & Demanded fan speed. \\
    24 & \texttt{PCNfR\_dmd}     & rpm     & Demanded corrected fan speed. \\
    25 & \texttt{W31}            & lbm/s   & HPT coolant bleed. \\
    26 & \texttt{W32}            & lbm/s   & LPT coolant bleed. \\
    \bottomrule
  \end{tabular}
\end{table}

The column grouping is meaningful: columns 1--2 are
identifiers, columns 3--5 are operating conditions, and
columns 6--26 are the 21 sensor measurements.

% -------------------------------------------------------------------
\section{Computing the Training RUL}
% -------------------------------------------------------------------

The training files do not contain a RUL column. It must be
derived from the data. The standard procedure is:

\begin{enumerate}[noitemsep]
  \item Group all rows by \texttt{unit\_number}.
  \item For each engine, find $t_{\max}$ --- the largest
        value of \texttt{time\_cycles} for that engine.
        This is the failure cycle.
  \item For every row of that engine:
        $\text{RUL} = t_{\max} - \texttt{time\_cycles}$
\end{enumerate}

This gives RUL $= 0$ at the failure cycle, counting upward
toward earlier cycles. Every engine in the training set ends
with RUL $= 0$.

\subsection{The piece-wise linear RUL cap}

A common preprocessing convention --- not specified in the
dataset or paper --- is to \textbf{cap the RUL} at a maximum
value (typically 125 cycles) for early cycles. The target
shape becomes piece-wise linear: flat at the cap value,
then linearly decreasing to zero.

Figure~\ref{fig:rulshape} shows the two shapes side by side.

\begin{figure}[H]
  \centering
  \includegraphics[width=0.88\textwidth]{11_rul_target_shape.png}
  \caption{Raw RUL (dashed blue) vs piece-wise linear capped
  RUL (solid teal, cap = 125). The cap removes the flat region
  at the top where precise long-range prediction is both
  impossible and uninformative.}
  \label{fig:rulshape}
\end{figure}

\begin{insightbox}
\textbf{Why cap the RUL?}
Early in an engine's life, the sensor signals carry little
degradation information --- the engine is healthy and its
behaviour is dominated by operating condition variability,
not wear. Asking a model to predict RUL = 180 precisely from
a healthy-looking signal is unreasonable and wastes model
capacity. The cap replaces this impossible task with a
flat ``healthy'' label, letting the model focus its capacity
on the degradation phase where prediction actually matters.

The cap value of 125 is a convention, not a physical constant.
Different papers use values from 100 to 150. When comparing
results across papers, always check whether a cap was applied
and what value was used.
\end{insightbox}

% -------------------------------------------------------------------
\section{The Multi-Condition Problem}
% -------------------------------------------------------------------

For FD002 and FD004, the six operating conditions cause a
problem that is not obvious until you look at the raw data.

\subsection{The problem}

Each of the six flight conditions produces a distinct range
of absolute sensor values. Fan inlet temperature $T_2$ at
sea level is hundreds of degrees Rankine higher than at
40,000\,ft. When an engine cycles through different flight
conditions from one cycle to the next, the raw sensor
readings jump between these bands --- not because of
degradation, but simply because of where the engine is
operating.

Figure~\ref{fig:condproblem} shows this effect clearly.

\begin{figure}[H]
  \centering
  \includegraphics[width=\textwidth]{12_multi_condition_problem.png}
  \caption{Top: raw $T_2$ values for one engine in FD002.
  The signal jumps between six bands due to condition
  switching --- not degradation. Bottom: the same signal
  after condition-wise normalisation. The condition bands
  collapse and a slow degradation drift becomes visible.}
  \label{fig:condproblem}
\end{figure}

Feeding raw sensor values from FD002 or FD004 into an ML
model trains the model to recognise operating conditions,
not degradation. The condition signal completely overwhelms
the health signal.

\subsection{The solution: condition-wise normalisation}

The standard fix is to normalize each sensor reading
relative to its expected healthy value \emph{at that
specific operating condition}. Figure~\ref{fig:normpipeline}
shows the full pipeline.

\begin{figure}[H]
  \centering
  \includegraphics[width=0.82\textwidth]{13_normalization_pipeline.png}
  \caption{The condition-wise normalisation pipeline for
  FD002 and FD004. Six clusters are identified from the
  three operating setting columns. Each sensor is
  normalised within its cluster using statistics computed
  from the healthy (early) cycles.}
  \label{fig:normpipeline}
\end{figure}

The procedure step by step:

\begin{enumerate}
  \item \textbf{Cluster} all cycles by operating condition
        using the three \texttt{op\_setting} columns.
        Six distinct clusters emerge naturally for
        FD002 and FD004.

  \item \textbf{Compute baseline statistics} for each
        sensor within each cluster, using only the early
        (healthy) cycles of training engines:
        \[
          \mu_{s,c} = \text{mean of sensor } s
                      \text{ in cluster } c
                      \text{ over healthy cycles}
        \]
        \[
          \sigma_{s,c} = \text{std of sensor } s
                         \text{ in cluster } c
                         \text{ over healthy cycles}
        \]

  \item \textbf{Normalise} every reading:
        \[
          x_{\text{norm}} = \frac{x - \mu_{s,c}}{\sigma_{s,c}}
        \]
\end{enumerate}

After this step, each sensor value expresses deviation from
the healthy baseline at that operating condition. Values near
zero mean healthy behaviour; drift away from zero indicates
degradation.

\begin{keybox}
\textbf{For FD001 and FD003} (single condition), this step
simplifies to a global normalisation --- one mean and
standard deviation per sensor across all training cycles.
The formula is the same; only one cluster is needed.
\end{keybox}

\begin{insightbox}
\textbf{What counts as ``healthy cycles''?}
There is no official definition. Common choices are:
the first $N$ cycles of each training engine (e.g.\ the
first 30), or all cycles of training engines with
RUL $\geq$ some threshold (e.g.\ RUL $\geq 125$).
The choice affects the baseline statistics and therefore
the normalized values. When comparing preprocessing
approaches across papers, this is a key parameter to check.
\end{insightbox}

% -------------------------------------------------------------------
\section{What the Data Does Not Contain}
% -------------------------------------------------------------------

It is worth being explicit about what is absent from the
files, because the simulation paper describes quantities that
never appear in the data:

\begin{center}
\begin{tabular}{@{}lp{7.5cm}@{}}
  \toprule
  Absent & Implication \\
  \midrule
  Health index $H(t)$              &
    Never provided. Must be inferred from sensors. \\
  Margin values (SmFan, SmHPC, \ldots) &
    Used only inside the simulation. Not in files. \\
  Hidden state $(e(t), f(t))$      &
    The actual degradation trajectory is never revealed. \\
  Fault mode label                 &
    FD003/FD004 do not indicate which module is degrading. \\
  Operating condition label        &
    The condition identity must be inferred from
    \texttt{op\_setting\_1/2/3}. \\
  RUL column in training files     &
    Must be computed: $t_{\max} - t$. \\
  \bottomrule
\end{tabular}
\end{center}

The algorithm sees only the 26 raw columns. Every other
quantity is part of the inference problem.

% -------------------------------------------------------------------
\section{Quick Reference}
% -------------------------------------------------------------------

\subsection{File naming}

\begin{center}
\begin{tabular}{@{}ll@{}}
  \toprule
  File & Purpose \\
  \midrule
  \texttt{train\_FD001.txt} & Training data, easiest sub-dataset \\
  \texttt{test\_FD001.txt}  & Test data, easiest sub-dataset \\
  \texttt{RUL\_FD001.txt}   & Test ground truth, easiest sub-dataset \\
  \texttt{train\_FD002.txt} & Training data, 6 conditions \\
  \texttt{test\_FD002.txt}  & Test data, 6 conditions \\
  \texttt{RUL\_FD002.txt}   & Test ground truth, 6 conditions \\
  \texttt{train\_FD003.txt} & Training data, 2 fault modes \\
  \texttt{test\_FD003.txt}  & Test data, 2 fault modes \\
  \texttt{RUL\_FD003.txt}   & Test ground truth, 2 fault modes \\
  \texttt{train\_FD004.txt} & Training data, hardest sub-dataset \\
  \texttt{test\_FD004.txt}  & Test data, hardest sub-dataset \\
  \texttt{RUL\_FD004.txt}   & Test ground truth, hardest sub-dataset \\
  \bottomrule
\end{tabular}
\end{center}

\subsection{Key numbers to remember}

\begin{center}
\begin{tabular}{@{}ll@{}}
  \toprule
  Quantity & Value \\
  \midrule
  Total files           & 12 \\
  Columns per row       & 26 \\
  Identifier columns    & 2 (unit, cycle) \\
  Operating settings    & 3 (cols 3--5) \\
  Sensor columns        & 21 (cols 6--26) \\
  Operating conditions  & 1 (FD001/FD003) or 6 (FD002/FD004) \\
  Fault modes           & 1 (FD001/FD002) or 2 (FD003/FD004) \\
  Common RUL cap        & 125 cycles (convention, not specified) \\
  \bottomrule
\end{tabular}
\end{center}

% -------------------------------------------------------------------
\end{document}
